%% kapa: exam=1 round=24 year=2013 period=2 form=A subject=accounting subject_ko=회계학 seq=9 number=9
%% source: Q-Net 감정평가사 기출문제 (cst00309) artlSeq=5073802 file=extracted/1008/5073802/2013년 제24회 감정평가사 1차시험 기출문제/2교시(A형).hwp
%% conversion: kordoc 4.24.51 -> tools/gen1.py (auto); review=proofread-sonnet; flags=none
\begin{question}{9}
(주)대한은 20×1년 1월 1일 공장건물을 ￦5,000,000에 취득하여 사용하기 시작하였다. (주)대한은 공장건물에 대해 취득 이후 원가모형을 적용하고 감가상각방법으로 정액법을 선택하였다(내용연수는 20년, 잔존가치는 ￦0). (주)대한은 20×6년 1월 1일 공장건물의 사용을 중단하고 투자부동산으로 전환하여 임대를 시작하였다. 분류변경 시점에서 잔여 내용연수는 10년으로 재추정되었고 공정가치모형을 적용하여 회계처리하기로 하였다. 20×6년 1월 1일 공장건물의 공정가치는 ￦3,800,000이었으며 20×5년 12월 31일까지 인식된 손상차손누계액은 ￦150,000이었다. 이러한 분류변경이 20×6년도 당기순이익에 미치는 영향은? (단, 20×6년도 건물의 공정가치는 변동이 없다.)
\begin{choices}
  \item ￦380,000 감소
  \item ￦180,000 감소
  \item ￦50,000 증가
  \item ￦150,000 증가
  \item ￦200,000 증가
\end{choices}
\end{question}
\answerall
\answernote{정답표에 1,2,3,4,5 모두 정답으로 표기(전원 정답)}
